\subsection{Make Straight His Paths: Confession}\label{sec:confession}

John the Baptist did not soften his preaching for the crowds who came out to
him:

\begin{quote}
He said therefore to the multitudes that went forth to be baptized by him: Ye
offspring of vipers, who hath shewed you to flee from the wrath to come? Bring
forth therefore fruits worthy of penance\dots\par
{\smallskip\noindent\small ---Luke 3:7--8\par}
\end{quote}

St Gregory the Great, in the homily the Office reads on the Fourth Sunday,
explains the warning in one sentence: ``\,`The wrath to come' in one sense
signifieth the great vengeance of the Latter Day: the sinner that repenteth not
of his sin now, will have no mean whereby to flee from punishment then.'' The
season that looks for the Judge also offers, before he comes, the way to meet
him without fear.

The way is the sacrament of penance. St Francis de Sales says it to the reader
of his \textit{Introduction to a Devout Life} as simply as it can be said:
``Our Saviour has left the holy sacrament of penance and confession to his
Church, that in it we might cleanse ourselves from all our iniquities, as often
as we should be defiled by them. Never suffer your heart then, O Philothea! to
remain long infected with sin, since you have so easy a remedy at hand''.

\subsubsection*{The land of confession}

When the priests and Levites came to ask John who he was, he confessed and did
not deny: he told the truth about himself. St Bernard, preaching on the Vigil
of Christmas, reminded his hearers that the Lord is born in Bethlehem of Juda,
and that the name of Juda means confession or praise.

The Council of Trent counts three acts of the penitent in the sacrament:
``contrition, confession, and satisfaction''. Of the three it puts contrition
first:

\begin{quote}
Contrition, which holds the first place amongst the aforesaid acts of the
penitent, is a sorrow of mind, and a detestation for sin committed, with the
purpose of not sinning for the future.\par
{\smallskip\noindent\small ---Council of Trent, Session 14\par}
\end{quote}

Sorrow, detestation, purpose: the act of contrition in
Appendix~\ref{sec:prayers} says all three. And the Council speaks of the fruit of
the sacrament: ``reconciliation with God, which sometimes, in persons who are
pious and who receive this sacrament with devotion, is wont to be followed by
peace and serenity of conscience, with exceeding consolation of spirit.''

\subsubsection*{An examination by the weeks of Advent}

These questions are drawn from the Gospels and Epistles of the season. Use them
with an examination on the commandments from a catechism or prayer book, not
instead of it. Take one part a night in the days before confession.

\begin{enumerate}
\item \textit{The first week: the works of darkness.} St Paul counts them:
rioting and drunkenness, impurity, contention and envy. Where have I made
provision for the flesh? Have I slept when I should have been awake to God,
with prayer neglected or put off?
\item \textit{The second week: the scandal and the reed.} Christ called
blessed the one who is not scandalized in him. Have I been ashamed of Christ, or
of his Cross? Have I bent like a reed to praise and blame, doing good to be seen
or leaving it undone for fear of what others would say? Have I passed by the
poor?
\item \textit{The third week: the truth about oneself.} John confessed what he
was and what he was not. Have I lied, flattered, boasted, or let a false
impression stand? Have I failed to see Christ standing among us, in the people
I live with?
\item \textit{The Ember Days: the handmaid.} Mary called herself the handmaid
of the Lord, and went with haste to her cousin. Where have I resisted God's will
when I knew it? Have I been slow to help those near me?
\item \textit{The fourth week: the crooked and the rough.} John cries that the
crooked ways must be made straight and the rough ways plain. Is there an
injustice in my dealings that I have not set right? Is there harshness in my
temper that I have excused? Have I judged others before the time, when only the
Lord can bring to light what is hidden in the heart?
\end{enumerate}

\subsubsection*{When to go}

Go in the third week of Advent, the week of the Ember Days, the Church's old
week of penance. If that cannot be, go before 17 December, when the last days
before Christmas begin. A confession made then leaves the rest of Advent to live
out the purpose of amendment. St Francis de Sales urged his reader to go far
more often than once a season, ``every week, and always, if possible, before you
communicate'', and he gave a reason that holds for anyone: ``by confession you
not only receive absolution from the venial sins you confess, but likewise
strength to avoid them, light to discern them well, and grace to repair all the
damage you may have sustained by them.''

Before the First Sunday, find out when confessions are heard in your parish in
Advent, and choose your day. Then go, and come back to the crib with nothing
between you and the Child.
